% Generated by Listas/make_statements.py from pratica-lecture-06.md, item 4. Do not edit.
\textbf{Source.} Muñoz-Garcia \& Toro-Gonzalez (2019), \emph{Strategy and Game Theory: Practice Exercises with Answers}, Ch. 4, Exercise 1 ``Ultimatum Bargaining Game'', p.~154 {[}PDF 166{]}. Reproduced verbatim (statement only; the book's worked answer follows it).



\begin{quote}

In the ultimatum bargaining game, a proposer is given a pie, normalized to a size \$1, and he is asked to make a monetary offer, \(x\), to the responder who, upon receiving the offer, only has the option to accept or reject it (as if he received an ``ultimatum'' from the proposer). If the offer \(x\) is accepted, then the responder receives it while the proposer keeps the remainder of the pie \(1 - x\). However, if he rejects it, both players receive a zero payoff. Operating by backward induction, the responder should accept any offer \(x\) from the proposer (even if it is low) since the alternative (reject the offer) yields an even lower payoff (zero). Anticipating such as a response, the proposer should then offer one cent (or the smallest monetary amount) to the responder, since by doing so the proposer guarantees acceptance and maximizes his own payoff. Therefore, according to the subgame perfect equilibrium prediction in the ultimatum bargaining game, the proposer should make a tiny offer (one cent or, if possible, an amount approaching zero), and the responder should accept it, since his alternative (reject the offer) would give him a zero payoff. However, in experimental tests of the ultimatum bargaining game, subjects who are assigned the role of proposer rarely make offers close to zero to the subject who plays as a responder.



Furthermore, sometimes subjects in the role of the responder reject positive offers, which seems to contradict our equilibrium predictions. In order to explain this dissonance between theory and experiments, many scholars have suggested that players' utility function is not as selfish as that specified in standard models (where players only care about the monetary payoff they receive). Instead, the utility function should also include social preferences, measured by the difference between the payoff a player obtains and that of his opponent, which gives rise to envy (when the monetary amount he receives is lower than that of his opponent) or guilt (when his monetary payoff is higher than his opponent's). In particular, suppose that the responder's utility is given by



\[u_R(x, y) = x + \alpha(x - y),\]



where \(x\) is the responder's monetary payoff, \(y\) is the proposer's monetary payoff, and \(\alpha\) is a positive constant. That is, the responder not only cares about how much money he receives, \(x\), but also about the payoff inequality that emerges at the end of the game, \(\alpha(x - y)\), which gives rise to either envy, if \(x < y\), or guilt, if \(x > y\). For simplicity, assume that the proposer is selfish, i.e., his utility function only considers his own monetary payoffs \(u_P(x, y) = y\), as in the basic model.



\textbf{Part (a)} Use a game tree to represent this game in its extensive form, writing the payoffs in terms of \(m\), the monetary offer of the proposer, and parameter \(\alpha\).



\textbf{Part (b)} Find the subgame perfect equilibrium. Describe how equilibrium payoffs are affected by changes in parameter \(\alpha\).



\textbf{Part (c)} Depict the equilibrium monetary amount that the proposer keeps, and the payoff that the responder receives, as a function of parameter \(\alpha\).

\end{quote}

